\chapter{Unicité du minimum}
\label{ch:convexite}

\begin{lecture}
La descente de gradient suppose un minimum unique. Sur une surface qui en compte
plusieurs, elle s'arrête au premier atteint et le résultat dépend du point de
départ. Le coût de la régression logistique n'en admet qu'un, ce qui se démontre
en vérifiant que sa courbure ne change jamais de signe.
\end{lecture}

\section{Courbure d'une fonction d'une variable}

La dérivée seconde mesure la variation de la pente. Si elle reste positive, la
pente croît sur tout l'intervalle et la concavité de la courbe reste tournée du
même côté. Une telle fonction est dite \emph{convexe} : elle ne peut pas présenter deux
minimums séparés par un maximum local, ce qui exigerait que la pente cesse de
croître.

\begin{figure}[htbp]
\centering
\begin{minipage}[t]{.48\textwidth}
\centering
\begin{tikzpicture}
\begin{axis}[width=.95\linewidth,height=4.8cm,axis lines=middle,
  title={convexe},xmin=-2.4,xmax=2.4,ymin=-.4,ymax=4,xtick=\empty,ytick=\empty,
  samples=120]
  \addplot[accent,very thick,domain=-2.2:2.2]{0.7*x^2+0.3};
  \addplot[only marks,mark=*,ink,mark size=2.4pt] coordinates {(0,0.3)};
\end{axis}
\end{tikzpicture}
\end{minipage}\hfill
\begin{minipage}[t]{.48\textwidth}
\centering
\begin{tikzpicture}
\begin{axis}[width=.95\linewidth,height=4.8cm,axis lines=middle,
  title={non convexe},xmin=-2.4,xmax=2.4,ymin=-.4,ymax=4,xtick=\empty,ytick=\empty,
  samples=160]
  \addplot[warning,very thick,domain=-2.2:2.2]{0.35*x^4-1.1*x^2+1.6};
  \addplot[only marks,mark=*,ink,mark size=2.4pt]
    coordinates {(-1.25,0.74) (1.25,0.74)};
\end{axis}
\end{tikzpicture}
\end{minipage}
\caption{À gauche, un minimum unique, atteint quel que soit le point de départ. À
droite, deux minimums locaux : la descente s'arrête dans celui vers lequel le
gradient initial la dirige, et ignore l'autre.}
\label{fig:convexe}
\end{figure}

\section{Courbure du coût logistique}

Reprenons la dérivée première (ch.~\ref{ch:gradient}) :
\[
  \frac{\partial J}{\partial w_j}
  =\frac{1}{m}\sum_{i=1}^{m}\bigl(\sigma(z_i)-y_i\bigr)x_{ij}.
\]
Dérivons une seconde fois, par rapport à un autre coefficient $w_k$ :

\begin{align}
  \frac{\partial^2 J}{\partial w_j\,\partial w_k}
  &= \frac{1}{m}\sum_{i=1}^{m}\frac{\partial}{\partial w_k}
     \bigl[\bigl(\sigma(z_i)-y_i\bigr)x_{ij}\bigr]
    &&\text{la dérivée entre dans la somme} \notag\\[2pt]
  &= \frac{1}{m}\sum_{i=1}^{m}x_{ij}\,\sigma'(z_i)\,
     \frac{\partial z_i}{\partial w_k}
    &&\text{chaîne, et } y_i \text{ est constant} \notag\\[2pt]
  &= \frac{1}{m}\sum_{i=1}^{m}x_{ij}\,\sigma(z_i)\bigl(1-\sigma(z_i)\bigr)x_{ik}
    &&\text{chapitre~\ref{ch:sigmoide}, et } \partial z_i/\partial w_k=x_{ik}
    \label{eq:hess}
\end{align}

Ces nombres, rangés dans un tableau carré dont la case $(j,k)$ porte
$\partial^2 J/\partial w_j\partial w_k$, forment la \emph{Hessienne} de $J$ : la
matrice de ses dérivées secondes. Posons
$s_i=\sigma(z_i)(1-\sigma(z_i))$, qui est le facteur commun à chaque terme.

\section{Positivité de la courbure}

Prenons une direction quelconque dans l'espace des coefficients, décrite par un
vecteur $\mathbf v=(v_0,\ldots,v_n)$. La courbure de $J$ le long de cette
direction vaut
\begin{align}
  \sum_{j}\sum_{k} v_j\,
  \frac{\partial^2 J}{\partial w_j\partial w_k}\,v_k
  &= \frac{1}{m}\sum_{i=1}^{m}s_i
     \left(\sum_j v_jx_{ij}\right)\left(\sum_k v_kx_{ik}\right)
    &&\text{on regroupe} \notag\\[2pt]
  &= \frac{1}{m}\sum_{i=1}^{m}s_i\left(\sum_j v_jx_{ij}\right)^{\!2}
    &&\text{les deux facteurs sont égaux} \notag\\[2pt]
  &\geq 0
\end{align}

La dernière ligne tient parce que chaque terme de la somme est le produit de deux
quantités positives ou nulles : un carré, toujours positif ou nul, et $s_i$, qui
vaut $\sigma(z_i)(1-\sigma(z_i))$ avec $\sigma$ strictement comprise entre $0$ et
$1$, donc lui aussi positif.

\begin{resultat}[Unicité du minimum]
La courbure de $J$ est positive ou nulle dans toutes les directions, quelle que
soit la position des coefficients. Le coût est donc convexe et n'admet pas de
minimum local secondaire. Tout point où le gradient s'annule est un minimum
global, et la descente y conduit quel que soit son point de départ : partir de
zéro relève de la reproductibilité, non de la qualité du résultat.
\end{resultat}

\section{Directions de courbure nulle}

Positif \emph{ou nul} : la nuance a des conséquences. La courbure s'annule dans
une direction $\mathbf v$ dès que $\sum_j v_jx_{ij}=0$ pour tous les élèves,
c'est-à-dire lorsqu'une combinaison des notes est identiquement nulle. L'ensemble
des minima est alors une droite, ou un sous-espace de dimension supérieure.

C'est exactement ce qui se produit avec deux colonnes proportionnelles. Si
Astronomy vaut $-100$ fois Defense Against the Dark Arts, augmenter un
coefficient et diminuer l'autre dans le bon rapport laisse tous les scores
inchangés : les prédictions restent identiques, mais les coefficients dérivent
sans se stabiliser.

\begin{vigilance}[Séparation complète]
Un second cas éloigne le minimum à l'infini. Si une droite sépare parfaitement
les deux groupes --- on dit alors que les données sont \emph{linéairement
séparables} --- multiplier tous les coefficients par $10$ rapproche encore les
probabilités de $0$ et de $1$, et fait donc encore baisser le coût. Chaque jeu de
coefficients en admet un meilleur et la suite diverge. La convexité établie
ci-dessus reste vraie : elle garantit l'unicité du minimum lorsqu'il existe, et
c'est son existence à distance finie que ce cas met en défaut
(annexe~\ref{ann:norme}).
\end{vigilance}
